\documentclass[11pt]{article}
\frenchspacing
\usepackage[colorlinks=false,pdfborder={0 0 1}]{hyperref}
\usepackage[
    top=0.65in,
    bottom=0.65in,
    left=0.7in,
    right=0.7in
]{geometry}

\setlength{\parindent}{0pt}
\setlength{\parskip}{0.35em}

\begin{document}

\noindent
{\small Madison T. Hobbs}
\hfill
{\small 22 September 2026}

\vspace{0.25em}

\begin{center}
    {\normalsize \textbf{16.S893 Proposal: AI-Assisted CFD Meshing and Verification for Hypersonic Simulations}}
\end{center}

\vspace{0.35em}
High-fidelity computational fluid dynamics (CFD) requires sufficient spatial resolution to capture the dominant flow gradients at a manageable cost. 
In hypersonic flow, those gradients are concentrated across shocks, boundary layers, surface-heating regions, and the wake. 
Mesh design depends on surface spacing, wall-normal resolution, boundary-layer growth, local shock and wake refinement, cell quality, and transitions in cell size. 
A mesh can satisfy geometric quality metrics while remaining too coarse to resolve the gradients that control pressure, heat flux, and wake structure. 
Uniform three-dimensional refinement rapidly increases cell count, memory demand, and node-hours. 
Efficient hypersonic meshing requires targeted resolution based on the local flow physics.

Mesh generation and verification are iterative and time-intensive. 
The usual workflow cycles through mesh construction, quality checks, solver execution, flowfield inspection, and targeted refinement. 
The resulting solution must also demonstrate grid convergence, with quantities of interest approaching stable values under systematic refinement. 
This is especially important in coupled fluid-material analyses, where CFD supplies the pressure and aerothermal loads that drive material response. 
Errors from inadequate fluid resolution propagate directly into the coupled prediction. 
Establishing a mesh-independent fluid solution can consume substantial time and compute before the primary multiphysics problem is addressed.

This project will assess whether an AI agent-assisted meshing workflow can reduce that burden while preserving standard CFD verification. 
Foundation models available through Parley will provide the reasoning layer. 
The test case will be the RAM-C II, a canonical flight-test benchmark for hypersonic blunt-body aerothermodynamics and wake-flow prediction \cite{Jones1972ElectrostaticProbePlasmaParameters}. 
The RAM-C II is well-suited to this study because its geometry, flight conditions, reference data, and existing numerical simulations are widely available and unrestricted. 
The demonstration will use the clean RAM-C II fluid problem to isolate the meshing task, with Pointwise through Glyph for grid generation and US3D, the University of Minnesota Computational Hypersonics Laboratory code for compressible and reacting-flow simulations, for the flow solution \cite{Candler2015US3DDevelopment}.
Pointwise Glyph scripts, agent instructions, verification tools, mesh-generation procedures, and US3D solutions will be shared in the class repository. The US3D source code will not be distributed.

The agent will begin with the RAM-C II geometry, freestream state, and solver setup, including the selected flow model, wall boundary condition, and angle of attack. 
The target quantities of interest will be surface pressure and heat flux, with computational constraints set by allowable cell count and solver cost. 
Reusable engineering skills will evaluate expected shock geometry, first-cell height for a target \(y^+\), wall-normal growth, surface-curvature resolution, cell-volume validity, skewness, orthogonality, aspect ratio, and local size transitions. 
Each skill will use deterministic calculations or literature-based criteria drawn from established hypersonic CFD and numerical methods.

Using these constraints, the agent will generate and modify the mesh by adjusting connector point counts, first-cell spacing, growth rates, and local shock, boundary-layer, and wake refinement. 
Topology changes will pause the workflow for human inspection in Pointwise and a dedicated topology review before modification. 
After each US3D solve, the agent will assess residual behavior, surface pressure, heat flux, and selected flowfield slices, then link observed deficiencies to specific mesh changes. 
It will record the affected region, parameter changed, and resulting solution response. 
The agent will then generate a topologically consistent coarse, medium, and fine grid family. 
Grid convergence in pressure and heat flux will be evaluated using solution differences, Richardson extrapolation, and Grid Convergence Index (GCI) \cite{Roache1994GridRefinement}.

The AI-assisted workflow will be compared with the conventional manual workflow using final solution quality, number of mesh iterations, failed solver runs, expert interventions, and total computational cost. 
Two failures of the AI-assisted workflow will be investigated. 
The first is a mesh that passes geometric quality checks but remains underresolved, with pressure, heat flux, or shock location still changing under refinement. 
The second is a mesh that is geometrically valid in Pointwise but produces unstable, poorly converged, or in the most extreme case, failed US3D runs. 

\newpage
\bibliographystyle{unsrt}
\bibliography{one_pager}
\end{document}